% 19_body.tex — Day4 산출물 4/5 (⑲) 교훈 등록부 (빈 템플릿 / 완성 예시 공용 본문)
% 드라이버: 19_교훈등록부_빈템플릿.tex (\wbblanktrue) / 19_교훈등록부_완성예시.tex (\wbblankfalse)
% 내용 출처: PM트랙/Day4_워크북.md §4.3(항목 가이드) §4.4(빈 템플릿) §4.5(온담 P1 완성 예시본)
\documentclass[10pt,a4paper]{article}
\usepackage[a4paper,margin=16mm,top=14mm,bottom=20mm]{geometry}
\def\DocNum{4}
\def\DocTitle{교훈 등록부}
\def\DocSub{⑲ Lessons Register (Lessons Report)}
\def\DocVariant{\B{빈 템플릿}{온담 P1 완성 예시}}
\usepackage{wbstyle}
\begin{document}
\landscapeon
\docheader
 {\Bg{[정식 명칭과 코드]}{차세대 주문·재고 통합 플랫폼 구축 (ONE ONDAM P1)}}
 {\Bg{[Lessons log 누적 n건 중 종료 정리 n건 / 발행일]}{Lessons log 2026-01 ~ 2027-05 누적 41건 중 종료 정리 15건 / 발행 2027-06-04}}
 {\Bg{[P1 PM]}{P1 PM (종료 보고서 부속)\rd{7mm}}}
 {\Bg{[수신 — 프로그램 PMO장·스폰서·내부감사]}{프로그램 PMO장 · 정미란 · 윤태호(사본)\rd{7mm}}}

% ------------------------------------------------------------------ 1 (가로) 교훈
\secbar{1. 교훈 — 상황 · 원인 · 교훈 · 적용 대상 · 소유자 · 증빙 \B{}{(적용 대상 약호: NP = 다음 프로젝트 양식 / STD = 사내 표준 / CTR = 계약 문안)}}
\guide{한 건에 여섯 칸. \textbf{상황은 숫자로}(12영업일·0.42억·+38분). \textbf{원인은 사람이 아니라 양식·규칙·구조}(“담당자가 놓쳤다”가 아니라 “양식에 그 칸이 없었다”). 교훈은 무엇을 바꾸면 되는지 한 문장. 적용 대상은 세 층 중 하나 이상 — 다음 프로젝트 양식(NP) / 사내 표준(STD) / 계약 문안(CTR) — 문서명과 칸까지. 소유자는 실명(PM은 곧 없어진다). 증빙은 어느 산출물의 어느 항목·행. 잘된 것도 같은 여섯 칸으로.}
{\scriptsize\setlength{\tabcolsep}{2pt}
\begin{longtable}{|C{11mm}|L{52mm}|L{44mm}|L{46mm}|L{40mm}|L{20mm}|L{26mm}|C{13mm}|}\hline
\hd{ID} & \hd{상황 (사실 · 일자 · 숫자)} & \hd{원인 (양식 · 규칙 · 구조)} & \hd{교훈 (한 문장)} & \hd{적용 대상 (NP / STD / CTR — 문서 · 칸)} & \hd{소유자 (실명)} & \hd{증빙 (문서 · 항목 · 행)} & \hd{유형} \\ \hline \endhead
\ifwbblank
\rd{12mm}L-01 & \g{[언제 무엇이 몇 — 숫자]} & \g{[어느 양식에 어느 칸이 없었나]} & \g{[무엇을 바꾸면 되나]} & \g{[STD ① 규정명 조항 / NP 양식 / CTR 계약 §]} & \g{[실명]} & \g{[⑫ 항목 n, 행 n]} & \g{[잘못 / 잘됨]} \\ \hline
\rd{12mm}L-02 & & & & & & & \\ \hline
\rd{12mm}L-03 & & & & & & & \\ \hline
\rd{12mm}L-04 & & & & & & & \\ \hline
\rd{12mm}L-05 & & & & & & & \\ \hline
\rd{12mm}L-06 & & & & & & & \\ \hline
\rd{12mm}L-07 & & & & & & & \\ \hline
\rd{12mm}L-08 & & & & & & & \\ \hline
\rd{12mm}L-09 & & & & & & & \\ \hline
\rd{12mm}L-10 & & & & & & & \\ \hline
\rd{12mm}L-11 & & & & & & & \\ \hline
\rd{12mm}L-12 & & & & & & & \\ \hline
\rd{12mm}L-13 & & & & & & & \\ \hline
\rd{12mm}L-14 & & & & & & & \\ \hline
\rd{12mm}L-15 & \g{[잘된 것 — 조건까지 적는다]} & & & & & & \g{[잘됨]} \\ \hline
\else
L-01 & 2026-04 성과보고 제4호 SPI 0.784를 “주의”로만 보고. 실제는 상용SW $-$6.0(조달 이연)·SI 실질 0 & 성과보고 양식에 원가 항목별 SV 분해 열이 없었음 & 지수는 열별로 분해해 보고하고, 지연의 종류(조달/개발/승인)를 적는다 & STD ① 성과보고 양식 항목 3 “항목별 SV·CV” 필수 / NP P2 & 프로그램 PMO장 & Day3 ⑫ 항목 9-1, 항목 3 & 잘못 \\ \hline
L-02 & R-05 실현(IS-02) — POS 화면 정의서 승인 대기 12영업일, 대체 결정권자 지정이 3주 늦음(06-03), 대기 비용 0.42억 & 헌장·RACI에 “결정권자 부재 시 대체자” 칸이 없어 리스크 대응(등록부 R-05 조치)이 실행 주체 없이 남음 & 결정 SLA가 있는 모든 A 역할에 대체 결정권자를 헌장에서 실명 지정한다 & STD ① 헌장 템플릿 §11 “대체 결정권자” 칸 / NP & 프로그램 PMO장 & ⑬ IS-02, ⑫ 신호 2 & 잘못 \\ \hline
L-03 & 5월 CV +2.70 급등을 “절감”으로 읽을 뻔함 — 라이선스 지급 조건 변경(50/50)의 인식 기준 차이 4.70 & 성과보고에 발생주의 조정 행이 없었음(8월 감리 준비 중 발견) & AC 인식 기준이 계약 변경으로 바뀌면 그 달부터 조정 행(발생주의 CPI)을 병기한다 & STD ① 성과보고 양식 항목 4 “조정 행” / NP & 프로그램 PMO장 & ⑫ 착시 1, 신호 3 & 잘못 \\ \hline
L-04 & 8월 결함 발견 주 40 → 12로 감소 — 품질 개선이 아니라 시험 인력 4명의 규제 검증 이동(시험 정지) & 결함표에 시험 노력(주당 실행 TC) 열이 없어 곡선만으로 판독 & 결함 곡선은 시험 노력과 반드시 병기하고, 평탄 구간은 “수렴/정지”를 판정한다 & STD ① 품질 통제 기록 양식 항목 2 “시험 노력” 열 / NP & 프로그램 PMO장 & ⑭ 항목 3 ③, ⑫ 신호 4 & 잘못 \\ \hline
L-05 & S18 3주 연장 결정(09-04) 시점에 A10 TF $-$10이 계산 가능했으나 예외 보고는 10-02 & 스프린트 연장 결정 양식에 “임계경로 TF 영향” 칸이 없었음 & 스프린트 연장·범위 재배치 결정에는 CPM TF 영향과 예외 보고 필요 여부를 같은 양식에서 판정한다 & NP 스프린트 결정 기록 양식 / STD ① & P1 PM → 프로그램 PMO장 & ⑫ 신호 5, 항목 8 & 잘못 \\ \hline
L-06 & 감리 1차 지적 AF-1-03·04가 발주사 미결정 원인인데 “수행사 설계 미흡”으로 기재 → 이의 서면·검수 보류 다툼·대기 비용 0.18 & 감리 계약에 지적 귀속 판정 기준(WBS 담당 열·RACI)이 없었음 & 감리 계약 부속서에 “귀속은 WBS 담당 열(발주/수행/공동)과 RACI를 기준으로 판정”을 사전 합의한다 & CTR 감리 계약 표준 부속서 3 / STD ③ & 강현도·정하늘 & ⑭ 항목 6, 이의 서면 09-29 & 잘못 \\ \hline
L-07 & 동시 지연 원칙(SCL 원칙 10)을 07-15 변경협의체에서야 합의 — 그 전 S10 판정은 발주사 단독 7일로만 처리 & SI 표준 계약 §14·§18에 여유 소유권·동시 지연 처리 문안이 없어 프로젝트 중 합의 & 여유 소유권(원칙 8)·동시 지연(원칙 10)을 계약 체결 시 문안에 넣는다 & CTR SI 표준 계약 §14·§18 / STD ③ & 강현도·최영주 & ⑮ 항목 6, Day2 ⑧ 항목 8 & 잘못 \\ \hline
L-08 & 대기 비용 판정(IS-02·AF 귀속)에서 수행사 대기 기록과 발주사 회의록이 일치한 덕에 판정이 1회로 끝남(07-15·10-14) & 주간 PM 회의록에 “결정 요청·응답 일자” 열을 착수 때부터 둔 것(동시 기록, contemporaneous records) & 결정 요청·응답·대기는 발생 당일 양쪽이 서명하는 동시 기록으로 남긴다 — 사후 재구성은 분쟁이 된다 & CTR §14 대기 기록 조항 / NP 주간 회의록 양식 & 강현도·P1 PM & ⑮ 항목 4 회신 4항, 항목 6 & \textbf{잘됨} \\ \hline
L-09 & D$-$2(02-24) API GW 운영 라이선스 발주가 CFO 직접 결재 대기(02-13~)로 발견 → 임시 라이선스 0.02억 & 내부통제 지침(5,000만 원 이상 CFO 결재, 2026-02-16 시행)에 프로젝트 창 앞 사전 결재 규정이 없고, G3 조달 상태 표가 “계약 내”로만 표기해 발주 단계 미추적 & 컷오버·게이트 앞 4주의 5,000만 원 이상 발주는 사전 결재 목록으로 일괄 결재하고, 조달 상태 표에 “발주·결재 단계” 열을 둔다 & STD ② 내부통제 지침 / STD ① 조달 상태 양식 & 재경팀장·정미란 / 강현도 & ⑯ 항목 9 & 잘못 \\ \hline
L-10 & 헌장 v1.1의 §5·§12 유보 표시가 G1(04-30) → 10-23까지 6개월 지속(CR-08·CR-10으로 해제) & 헌장 템플릿에 유보 항목의 해제 기한·책임자 칸이 없음 & 헌장 유보 항목은 다음 게이트까지 해제하거나 CCB에 상정한다 & STD ① 헌장 템플릿 “유보 항목·해제 기한” & 프로그램 PMO장 & ⑪ 항목 9, Day2 §6 & 잘못 \\ \hline
L-11 & Rules of Credit(EV 측정 규칙)이 원가 기준선 부속으로 정식 첨부된 것이 성과보고 제9호(10월) — G1 기준선에는 PV 배분 규칙만 있어 측정 시차 착시(1.96) & 원가 기준선 템플릿에 RoC 부속이 없음 & PV 배분 규칙과 EV 인정 규칙을 기준선 승인 시 한 문서로 서명한다 & STD ① 원가 기준선 템플릿 부속 2 & 프로그램 PMO장 & ⑫ 항목 2, 착시 3 & 잘못 \\ \hline
L-12 & 컨틴전시 R-03 0.50 집행(08-12 결재)의 8월 CCB 보고 누락 → 감리 지적 AF-1-08 → 10-15 소급 & 집행 결재와 CCB 안건이 별개 절차 & 컨틴전시 집행 결재 시 다음 CCB 안건이 자동 생성되도록 절차를 묶는다(11월부터 적용, 이후 누락 0) & STD ① 프로젝트 관리 규정 컨틴전시 조항 & 프로그램 PMO장 & ⑭ AF-1-08, ⑬ 항목 3 & 잘못→고침 \\ \hline
L-13 & 프리모템 원문 61건 중 15건만 리스크로 승격 — 승격되지 않은 6번(개발자 이탈)·12번(보안)·13번(교육 이탈)이 IS-11·IS-08·IS-04로 실현 & 승격 기준이 득표순이라 “발생 시 대응 시간이 짧은 것”이 탈락 & 프리모템 승격 기준에 득표 외 “대응 리드타임”을 추가하고, 미승격 원문은 감시 목록으로 분기 재검토 & STD ① 리스크 관리 계획 템플릿 §2 승격 기준 / NP & 프로그램 PMO장 & ⑬ 항목 2 출처 분포, Day1 ⑤ & 잘못 \\ \hline
L-14 & P1-02 포인트 이중차감(D+5) — 개정법 ‘훼손’ 해당 여부를 단정할 수 없는 상태에서 49h 40m 내 신고 결정 & 운영 문서에 개인정보 사건 판단 절차(인지 시각 기록·CPO 의견 시한·통지/신고 분리)가 없어 결정 회의를 소집해 처리 & 규제 적용은 단정하지 않되 시한은 단정한다 — 인지 기록 → CPO 24h 의견 → 통지 즉시 → 신고 판단 48h & STD ② 개인정보 내부관리계획(최영주) / NP 운영 문서 8종 & 최영주·박준영 & ⑰ 항목 7 & 잘됨(절차화 필요) \\ \hline
L-15 & 두 번 컷오버(02-26·03-27)가 모두 정시·트리거 발동 0 — 조건: 1차 하이퍼케어 지표(심각도 2 이상 ≤ 5)를 2차 go 조건으로, 리허설 1회차를 전체 볼륨으로 실행해 창 초과(41h)를 미리 관측 & 런북을 2회 분리 작성하고 R-14 트리거를 게이트 조건으로 연결(G1 의결 → Day3 R-14 → ⑯) & 분리 컷오버는 앞 컷오버의 안정화 지표를 뒤 컷오버의 게이트 조건으로 묶고, 리허설 1회차는 실패를 보기 위해 돌린다 & STD ① 프로그램 컷오버 표준(런북 양식·게이트 3·PoNR·트리거) / NP P2·P5 & 프로그램 PMO장·박준영 & ⑯ 항목 3·6·7 & \textbf{잘됨} \\ \hline
\fi
\end{longtable}}
\B{\small\g{원천 분포(조기 신호 / 착시 / 감리·계약 / 이슈·리스크 / 절차·거버넌스 / 잘된 것) — 어느 워크북 산출물에서 왔는가: }\hrulefill\par}{\small \textbf{원천 분포} — 조기 신호 5(⑫ 항목 9) / 착시(⑫ 항목 7) 1 / 감리·계약(⑭·⑮) 3 / 이슈·리스크(⑬) 2 / 절차·거버넌스(⑯·⑪·Day1) 3 / 잘된 것 1.\par}

% ------------------------------------------------------------------ 2 (가로) 사내 표준 개정 요청
\secbar[80mm]{2. 사내 표준 개정 요청 \B{}{(3건, 06-04 발송)}}
\guide{교훈을 적용 대상별로 묶어 개정 요청서로 만든다 — \textbf{개정 요청은 교훈 등록부의 부속이 아니라 결과물이다.} 표준명(소유 부서)·개정 조항·요청 문안(조항 번호와 새 문장)·관련 교훈·요청자 → 수신(표준의 소유 부서)·기한·상태. “검토 바람”은 요청이 아니다. 교훈 등록부를 아카이브에 넣고 끝내면 다음 프로젝트에서 같은 새벽에 다시 발견된다.}
{\footnotesize\setlength{\tabcolsep}{2.5pt}
\begin{tabularx}{\linewidth}{|C{6mm}|L{36mm}|Y|L{30mm}|L{34mm}|L{34mm}|L{18mm}|}\hline
\hd{\#} & \hd{표준 (소유)} & \hd{개정 조항 · 요청 문안} & \hd{관련 교훈} & \hd{요청자 → 수신} & \hd{기한} & \hd{상태} \\ \hline
\ifwbblank
\rd{18mm}1 & \g{[프로젝트 관리 규정·양식집 (PMO)]} & \g{[① 양식명: 바뀌는 칸 ② … — 조항 번호와 새 문장]} & \g{[L-nn·nn]} & \g{[실명 → 실명·직책]} & \g{[의결 회의체·일자, 적용 시점]} & \g{[접수 일자]} \\ \hline
\rd{18mm}2 & \g{[내부통제 지침 (재무)]} & \g{[조항에 신설할 단서 — 문장으로]} & & & & \\ \hline
\rd{18mm}3 & \g{[구매 계약 표준 문안 (구매·법무)]} & \g{[계약 § — 문안]} & & & & \\ \hline
\else
1 & \textbf{프로젝트 관리 규정·양식집}(프로그램 PMO) & ① 성과보고 양식: 항목별 SV·CV 분해 열, 발생주의 조정 행 ② 헌장 템플릿: 대체 결정권자 칸, 유보 항목 해제 기한 ③ 원가 기준선 템플릿: RoC 부속 2 ④ 품질 기록: 시험 노력 열 ⑤ 컨틴전시 집행 → CCB 안건 자동 생성 ⑥ 리스크 관리 계획: 승격 기준에 대응 리드타임 ⑦ 프로그램 컷오버 표준(런북·게이트·PoNR·트리거·분리 컷오버 연결) ⑧ BC 편익 표: 분모·측정 책임자 실명 필수(⑳ 항목 1) & L-01·02·03·04·05 · L-10·11·12·13·15 & P1 PM → 프로그램 PMO장 & 2027-07 운영위 의결, P2 G2 전 적용 & 접수 06-08 \\ \hline
2 & \textbf{재무본부 내부통제 지침}(재경팀장·정미란) & 5,000만 원 이상 CFO 직접 결재 조항에 단서 신설: “프로그램 게이트·컷오버 창 앞 4주 이내 발주 예정 건은 소관 PM이 사전 결재 목록을 제출하고 CFO가 일괄 결재한다. 결재 대기 5영업일 초과 건은 재경팀이 PM에 통보한다” + 개인정보 내부관리계획에 사건 판단 시한 절차 & L-09, L-14 & P1 PM → 재경팀장·최영주, 정미란 & 2027-08 지침 개정 & 접수 06-08 \\ \hline
3 & \textbf{구매 계약 표준 문안}(강현도·최영주) & ① SI 표준 계약 §14: 여유 소유권(SCL 원칙 8)·동시 지연(원칙 10)·대기 기록(동시 기록·양측 서명) 문안 ② 감리 계약 부속서 3: 지적 귀속 판정 기준 = WBS 담당 열·RACI ③ 조달 상태 양식: 발주·결재 단계 열 & L-06·07·08·09 & P1 PM → 강현도·최영주 & 2027-09 P2 후속 계약부터 & 접수 06-08 \\ \hline
\fi
\end{tabularx}}

% ------------------------------------------------------------------ 3 (가로) 다음 프로젝트 착수 패키지 양식
\secbar[100mm]{3. 다음 프로젝트 착수 패키지에 들어갈 양식 \B{}{(P2 G2 · P5 착수 · 2028 후속 과제 공통)}}
\guide{양식 · 바뀐 칸 · 관련 교훈. 교훈이 “다음 프로젝트의 착수 패키지”라는 물리적 형태로 남을 때만 조직에 남는다. 정리에서 제외한 교훈(운영 개선 등록부·타 프로젝트 소관·재발 가능성 낮음)의 분류와 건수를 아래에 적는다.}
{\footnotesize\setlength{\tabcolsep}{2.5pt}
\begin{tabularx}{\linewidth}{|L{50mm}|Y|L{40mm}|}\hline
\hd{양식} & \hd{바뀐 칸} & \hd{교훈} \\ \hline
\ifwbblank
\rd{6mm}\g{[헌장 템플릿 v2]} & \g{[§ 칸]} & \g{[L-nn]} \\ \hline
\rd{6mm}\g{[성과보고 양식 v2]} & & \\ \hline
\rd{6mm}\g{[원가 기준선 템플릿 v2]} & & \\ \hline
\rd{6mm}\g{[품질 통제 기록 v2]} & & \\ \hline
\rd{6mm}\g{[리스크 관리 계획 v2]} & & \\ \hline
\rd{6mm}\g{[컨틴전시 집행 결재서 v2]} & & \\ \hline
\rd{6mm}\g{[컷오버 런북·SAC 양식 v1]} & & \\ \hline
\rd{6mm}\g{[주간 PM 회의록 v2]} & & \\ \hline
\rd{6mm}\g{[BC 편익 표 v2]} & & \\ \hline
\else
헌장 템플릿 v2 & §11 대체 결정권자(실명) / 유보 항목·해제 기한 & L-02·L-10 \\ \hline
성과보고 양식 v2 & 항목 3 항목별 SV·CV / 항목 4 조정 행 / 항목 8 스프린트 연장 TF 영향 & L-01·L-03·L-05 \\ \hline
원가 기준선 템플릿 v2 & 부속 2 Rules of Credit & L-11 \\ \hline
품질 통제 기록 v2 & 항목 2 시험 노력 열 & L-04 \\ \hline
리스크 관리 계획 v2 & 승격 기준·감시 목록 & L-13 \\ \hline
컨틴전시 집행 결재서 v2 & CCB 안건 자동 생성 & L-12 \\ \hline
컷오버 런북·SAC 양식 v1 & ⑯ 전체(이 워크북 §1.4) & L-15 \\ \hline
주간 PM 회의록 v2 & 결정 요청·응답·대기 열(양측 서명) & L-08 \\ \hline
BC 편익 표 v2 & 분모·측정 책임자 실명·분모 확장 조건 & ⑳ \\ \hline
\fi
\end{tabularx}}
\B{\small\g{Lessons log n건 중 정리 제외 n건의 분류(운영 개선 등록부 n / 타 프로젝트 소관 n / 재발 가능성 낮음 n) — 목록은 아카이브: }\hrulefill\par}{\small Lessons log 41건 중 정리에서 제외한 26건은 운영 개선 등록부(CIR, 박준영)로 12건, P3·P4 소관으로 9건, 재발 가능성 낮음으로 5건 분류(목록은 아카이브).\par}
\vspace{3mm}
\noindent{\footnotesize
\begin{tabularx}{\linewidth}{|Y|Y|Y|Y|}\hline
\hd{작성 — P1 PM (서명 · 일자)} & \hd{접수 — 프로그램 PMO장 (서명 · 일자)} & \hd{확인 — 스폰서 (서명 · 일자)} & \hd{사본 — 내부감사 (수신 일자)} \\ \hline
\Bg{[P1 PM]}{P1 PM 2027-06-04\rd{10mm}} & \Bg{[프로그램 PMO장]}{프로그램 PMO장 06-08} & \Bg{[스폰서]}{정미란 06-08} & \Bg{[내부감사팀장]}{윤태호 06-08} \\ \hline
\end{tabularx}}
\landscapeoff
\end{document}
